% ============================================================
% Introduction — in astronomy this section ALSO carries the
% literature review (there is no separate Related Work). Runs
% longer than in CS venues (2–5 pages), funnels from the field
% to this paper, with citations woven in via \citet (textual)
% and \citep (parenthetical). See references/writing-guide.md.
% ============================================================
\section{Introduction}
\label{sec:intro}

% ---- Move 1: Astrophysical context (1–2 paragraphs) ----
% Open with the big picture and why it matters physically.
% \citet{1926ApJ....64..321H} is textual ("Hubble 1926 showed"),
% \citep{1980ApJ...236..351D} is parenthetical "(Dressler 1980)"
% — choosing between them is the main citation-craft in astro.
Since the pioneering work of \citet{1926ApJ....64..321H}, this field has provided fundamental clues to our understanding of the problem.

% ---- Moves 2–3: Prior work and the open question (3–6 paragraphs) ----
% Unlike CS papers, this is where related work lives. Organize
% by theme/progression of ideas, not paper-by-paper. Each
% paragraph: what was done (cited), what it established, what
% it left open. Build the tension that the paper resolves.
% End with the specific gap stated in one crisp sentence.

% ---- Move 4: This paper (1–2 paragraphs) ----
% "In this paper, we ..." — what we do, how, and a preview of
% the headline result with numbers. Astronomy papers close the
% intro with a compact prose summary of findings rather than a
% bulleted contribution list; the final paragraph may outline
% the paper structure ("Section 2 describes the data ...") —
% optional in modern practice, keep it only if the structure is
% unusual.
